\documentclass[11pt]{article}
\usepackage[utf8]{inputenc}
\usepackage[T1]{fontenc}
\usepackage{lmodern}
\usepackage[margin=0.9in]{geometry}
\usepackage{amsmath,amssymb,amsthm}
\usepackage{booktabs}
\usepackage{longtable}
\usepackage{array}
\usepackage{enumitem}
\usepackage[hidelinks]{hyperref}
\setlength{\emergencystretch}{3em}

\newtheorem{theorem}{Theorem}
\newtheorem{proposition}[theorem]{Proposition}
\newtheorem{corollary}[theorem]{Corollary}
\theoremstyle{definition}
\newtheorem{definition}[theorem]{Definition}
\newtheorem{example}[theorem]{Example}
\newtheorem{remark}[theorem]{Remark}

\newcommand{\G}{\mathcal G}
\newcommand{\C}{\mathcal C}
\newcommand{\D}{\mathcal D}
\newcommand{\Nat}{\operatorname{Nat}}
\newcommand{\Aut}{\operatorname{Aut}}
\newcommand{\Fix}{\operatorname{Fix}}
\newcommand{\Img}{\operatorname{im}}
\newcommand{\Sel}{\operatorname{Sel}}
\newcommand{\EqI}{\mathrel{\sim_I}}

\title{Natural Selection Under Symmetry:\\
Fixed-Point Obstructions, Finite Countermodels, and Minimum Extensions}
\author{The Aether-Flow Research Project---draft/control}
\date{23 July 2026}

\begin{document}
\maketitle

\begin{center}
\fbox{\begin{minipage}{0.93\linewidth}
\textbf{Status and authority.}
This manuscript is a task-local \texttt{draft/control},
\texttt{proposal-only} science artifact.  It integrates previously registered
mathematical results without adopting a source category, selector law,
quotient, physical gauge, or ontology extension.  It does not establish
general source equivalence (\(\mathrm{EqSrc}\)), a global no-go theorem,
future-extension impossibility, proof authority, peer review, submission
eligibility, publication authority, benchmark promotion, or a completed
derivation of general relativity.
\end{minipage}}
\end{center}

\begin{abstract}
Let a small groupoid act functorially on sets of eligible choices.  Natural
deterministic selectors are classified by a product of automorphism-fixed
loci, which gives exact existence, uniqueness, and nonuniqueness criteria.
A natural map from choices to relation data supplies a separate criterion for
uniqueness of the induced relation, so representative nonuniqueness need not
imply relational nonuniqueness.  Finite sign, line, root, grading, and
relation-image models exhibit the alternatives.  We then organize candidate
repairs in a typed eight-atom declaration envelope and prove that, inside any
finite sound obstruction-clause model, every successful repair contains an
inclusion-minimal hitting signature.  Finally, we compare representative
selection with conditional quotient-first constructions and expose their
source-derivation, interface-descent, and physical-interpretation burdens.
The results are exact on their declared mathematical domains.  They do not
derive those domains from the current Aether-Flow ontology or establish a
universal physical obstruction.
\end{abstract}

\section{Question and contribution}

Many structural models present a family of mathematically equivalent-looking
descriptions and then require either one distinguished representative or a
representative-free object.  The basic question is:
\begin{quote}
When can a choice be made coherently under every admitted structural
isomorphism, when is that choice unique, and what additional structure is
unavoidably declared when the answer is negative?
\end{quote}

The first result below is stated without project-specific vocabulary.  The
later sections instantiate it on finite source models, relate failures to a
typed extension envelope, and compare selection with a quotient route.  The
individual category-theoretic ingredients are not claimed here as new.  This
internal draft is a project-specific, source-traced synthesis of the
fixed-locus classification, finite countermodels,
minimum-extension logic, and conditional quotient implications, with their
claim boundaries made explicit.

\section{Natural selectors on a groupoid}

\subsection{Data}

Let \(\G\) be a small groupoid and let
\[
  S:\G\longrightarrow\mathbf{Set}
\]
be a functor.  Write \(I=\pi_0(\G)\) for the set of connected components.
For each \(i\in I\), supply a representative object \(r_i\).  For every
object \(X\) in component \(i=[X]\), supply an isomorphism
\[
  \tau_X:r_i\longrightarrow X,\qquad \tau_{r_i}=1_{r_i}.
\]
Define
\[
  H_i:=\Aut_{\G}(r_i),\qquad
  F_i:=S(r_i)^{H_i}
  =\{x\in S(r_i):S(h)x=x\ \text{for every }h\in H_i\}.
\]
A deterministic natural selector is a natural transformation
\[
  \sigma:1_{\G}\Rightarrow S,
\]
where \(1_{\G}\) is the constant singleton functor.  Let
\(\Sel(\G,S):=\Nat(1_{\G},S)\).

\begin{theorem}[Fixed-locus classification]
\label{thm:classification}
Evaluation at the supplied representatives is a bijection
\[
  \Sel(\G,S)\xrightarrow{\;\cong\;}\prod_{i\in I}F_i,\qquad
  \sigma\longmapsto(\sigma_{r_i})_{i\in I}.
\]
The inverse sends \(x=(x_i)_{i\in I}\) to
\[
  \sigma^x_X=S(\tau_X)(x_{[X]}).
\]
\end{theorem}

\begin{proof}
If \(\sigma\) is natural and \(h\in H_i\), naturality at
\(h:r_i\to r_i\) gives \(S(h)(\sigma_{r_i})=\sigma_{r_i}\); evaluation
therefore lands in the stated product.

Conversely, take \(x\in\prod_iF_i\).  For an isomorphism \(f:X\to Y\) in
component \(i\), set
\[
  a:=\tau_Y^{-1}f\tau_X\in H_i.
\]
Because \(x_i\) is fixed by \(H_i\),
\[
  S(f)\sigma^x_X
  =S(f\tau_X)x_i
  =S(\tau_Ya)x_i
  =S(\tau_Y)x_i
  =\sigma^x_Y.
\]
Thus \(\sigma^x\) is natural.  Its value at \(r_i\) is \(x_i\).
Conversely, naturality of \(\sigma\) at \(\tau_X\) gives
\(\sigma_X=S(\tau_X)(\sigma_{r_i})\), so evaluation and reconstruction are
inverse.
\end{proof}

\begin{corollary}[Existence and uniqueness]
\label{cor:existence}
The following statements hold.
\begin{enumerate}[label=(\alph*)]
  \item \(\Sel(\G,S)=\varnothing\) exactly when
        \(\prod_iF_i=\varnothing\).
  \item If some \(F_i\) is empty, then no selector exists.
  \item If every \(F_i\) is nonempty, a selector exists exactly when a product
        element is available.  This is automatic for finite \(I\), for a
        supplied tuple, or under a declared suitable choice principle.
  \item Exactly one selector exists if and only if every \(F_i\) is a
        singleton.
  \item At least two selectors exist if and only if some \(F_i\) contains two
        points and the complementary product
        \(\prod_{j\ne i}F_j\) is nonempty.
\end{enumerate}
\end{corollary}

\begin{proof}
Apply elementary product-set facts to
Theorem~\ref{thm:classification}.  The qualification in (c) is essential:
in ZF, factorwise nonemptiness of an arbitrary family does not itself supply
a product element.
\end{proof}

\begin{remark}[What the theorem does not say]
The automorphisms are the arrows declared structural in \(\G\).  The theorem
does not show that they are physical gauge transformations.  A fixed point is
a coherent mathematical choice relative to the supplied functor; it is not
thereby source-derived, physically preferred, operationally measurable, or
empirically adequate.
\end{remark}

\subsection{Relation images are a separate question}

Let \(E:\G\to\mathbf{Set}\) be another functor and let
\(K:S\Rightarrow E\) be a natural transformation.  Composition defines
\[
  K_*:\Sel(\G,S)\longrightarrow\Nat(1_{\G},E),\qquad
  K_*(\sigma)=K\circ\sigma .
\]
Write \(R_i=K_{r_i}(F_i)\).

\begin{theorem}[Relation-image uniqueness]
\label{thm:relation-image}
The map \(K_*\) is well defined.  If \(\Sel(\G,S)\ne\varnothing\), then
\[
  |\Img K_*|=1
  \quad\Longleftrightarrow\quad
  |R_i|=1\ \text{for every }i\in I.
\]
Thus multiple selectors can induce one relation section, while two fixed
choices with distinct \(K\)-images produce distinct relation sections when
the complementary product is nonempty.
\end{theorem}

\begin{proof}
Naturality of \(K\) makes \(K\circ\sigma\) natural.  Under
Theorem~\ref{thm:classification}, the representative values of
\(K_*(\sigma)\) are exactly \(K_{r_i}(x_i)\).  If every \(R_i\) is a
singleton, all selector tuples have the same relation values.  Conversely,
if some \(R_i\) contains two values, combine their preimages with one tuple
from the nonempty complementary product to obtain two selector tuples whose
relation sections differ at \(r_i\).
\end{proof}

\subsection{Noninvertible arrows add equalizer obligations}

Let \(\C\) be a small category with wide core \(\G=\operatorname{Core}(\C)\),
and let \(S:\C\to\mathbf{Set}\).

\begin{theorem}[Full-category equalizer guard]
\label{prop:equalizer}
Restriction gives an injection
\[
  \Nat(1_{\C},S)\hookrightarrow\Nat(1_{\G},S).
\]
Its image consists exactly of the core selectors satisfying
\[
  S(f)(\sigma_X)=\sigma_Y
\]
for every noninvertible arrow \(f:X\to Y\) of \(\C\).
\end{theorem}

\begin{proof}
The wide core has every object, so two natural transformations agreeing after
restriction have identical components.  A core selector extends to
\(\C\) precisely when its components satisfy the remaining naturality
equations.
\end{proof}

The fixed-locus product therefore classifies the groupoid core.  It is not a
substitute for the additional equalizer checks on a larger category.

\section{Finite countermodels and positive controls}

The companion JSON file reproduces the following finite calculations.  They
are examples on declared mathematical actions, not physical systems.

\begin{longtable}{>{\raggedright\arraybackslash}p{0.25\linewidth}
                  >{\raggedright\arraybackslash}p{0.18\linewidth}
                  >{\raggedright\arraybackslash}p{0.47\linewidth}}
\toprule
Model & Exact fixed-locus result & Consequence \\
\midrule
\endhead
Sign swap & \(C_2\) acts transitively on two signs; fixed set empty
& No selector on that declared one-object groupoid. \\
Tagged singleton & trivial action on one supplied tag
& One conditional selector; tag provenance is an added assumption. \\
Three lines in \(\mathbb F_2^2\) & \(\mathrm{GL}(2,2)\) has order \(6\),
acts transitively on three lines; fixed set empty
& No invariant line selector under the full declared action. \\
Marked line & its stabilizer has order \(2\) and fixes exactly the mark
& One conditional selector after adding line data. \\
Two-cycle root & \(C_2\) swaps two roots; fixed set empty
& No root selector on the unmarked cycle. \\
Four-cycle root & \(C_4\) acts transitively on four roots; fixed set empty
& No root selector; the four candidates nevertheless fall into two declared
parity relation classes. \\
Loaded root & root-preserving choice space is a singleton
& One conditional selector; root provenance remains open. \\
Graded four-chain & order automorphisms are trivial and grade zero is unique
& One conditional minimum; the grading is supplied structure. \\
Trivial two-choice control & two fixed points
& Two selectors.  An injective \(K\) gives two relation images while a
constant \(K\) gives one. \\
\bottomrule
\end{longtable}

These controls isolate three logically different outcomes: an empty fixed
locus obstructs a selector; a singleton fixed locus gives a conditional
unique selector; and a multiple fixed locus leaves representative
underdetermination.  The last row directly prevents the invalid inference
``multiple selectors imply multiple relations.''
For the empty-fixed-locus rows, pointwise \(K\)-values may still be available
on eligible choices, but they are not elements of \(\operatorname{Img}K_*\):
the domain of \(K_*\) is empty because no natural selector exists.  The
companion data therefore records eligible relation values separately from
induced natural relation sections.

\section{A typed extension envelope}

Suppose a failed selector construction is repaired by declaring additional
source-side structure.  Use the eight bookkeeping atoms
\[
  \mathcal A=\{R,O,Q,G,M,I,B,D\},
\]
standing respectively for root/partition data, orientation/line data,
quotient/discriminator data, grade/order data, measure/kernel data,
initial-state data, boundary data, and transition/evolution data.

\begin{proposition}[Exact declaration envelope]
\label{prop:b8}
The syntactic declaration envelope
\[
  \mathcal L_{\mathrm{decl}}=(\mathcal P(\mathcal A),\subseteq)
\]
is the Boolean lattice \(B_8\).  It has \(256\) nodes, \(1024\) directed
cover edges, and rank sizes
\[
  (1,8,28,56,70,56,28,8,1).
\]
\end{proposition}

\begin{proof}
There are \(2^8=256\) subsets.  A cover adds one absent atom, so the number
of directed covers is
\(\sum_{E\subseteq\mathcal A}(8-|E|)=8\,2^7=1024\).
Rank \(k\) has \(\binom{8}{k}\) nodes.
\end{proof}

Boolean inclusion is only a declaration check.  Scientific assumption
strength requires a well-typed reduct between realized extensions over the
same baseline and target.  The admissible extension poset need not be all of
\(B_8\), and the Boolean rank is neither a plausibility score nor an adoption
rank.

For a finite structural automorphism group \(K_X\) acting on eligible choices
\(\mathcal D(X)\), let an extension package \(e_E\) have stabilizer
\[
  H_E=\bigcap_{a\in E}\operatorname{Stab}_{K_X}(e_a)
\]
and an \(H_E\)-stable compatible-choice set \(C_E(e_E)\).  Define
\[
  \Sigma(E,e_E)=C_E(e_E)\cap\Fix_{\mathcal D(X)}(H_E).
\]

\begin{proposition}[Compatible fixed-locus trichotomy]
On the realized one-object finite groupoid,
\[
\begin{array}{rcl}
|\Sigma|=0 &\Longleftrightarrow& \text{no compatible deterministic selector},\\
|\Sigma|=1 &\Longleftrightarrow& \text{one compatible deterministic selector},\\
|\Sigma|>1 &\Longleftrightarrow& \text{residual deterministic underdetermination}.
\end{array}
\]
\end{proposition}

\begin{proof}
A compatible natural choice is exactly an element satisfying both
compatibility and stabilizer-fixedness.
\end{proof}

Uniqueness is not monotone under declaration inclusion: a larger package may
reduce the stabilizer and thereby enlarge the raw fixed locus.  Compatibility
and typing must do substantive work.

\section{Finite minimum extensions}

Let \(U\) be a finite declared universe of extension capabilities.  Let
\(J=\{1,\dots,n\}\) index obstruction burdens, with each nonempty
\(C_j\subseteq U\) listing capabilities that can discharge burden \(j\).
For \(E\subseteq U\), define
\[
  \operatorname{Esc}(E)
  \quad\Longleftrightarrow\quad
  E\cap C_j\ne\varnothing\ \text{for every }j\in J.
\]

\begin{theorem}[Finite minimum-extension signatures]
\label{thm:hitting}
For finite \(U\) and nonempty clauses \(C_j\):
\begin{enumerate}[label=(\alph*)]
  \item \(\operatorname{Esc}(E)\) holds exactly when \(E\) is a hitting set
        of \(\{C_j:j\in J\}\);
  \item every escaping \(E\) contains an inclusion-minimal escaping subset;
  \item the collection
        \[
          \mathcal H_{\min}
          =\{H\subseteq U:\operatorname{Esc}(H)
          \text{ and no proper subset of }H\text{ escapes}\}
        \]
        is finite and nonempty; and
  \item every successful construction represented in a sound clause model
        realizes at least one \(H\in\mathcal H_{\min}\).
\end{enumerate}
\end{theorem}

\begin{proof}
Part (a) is the definition of a hitting set.  Starting from a finite escaping
\(E\), delete any dispensable element while the escape predicate remains
true.  Finite descent terminates at an inclusion-minimal escaping subset,
proving (b).  The powerset of \(U\) is finite and \(U\) hits every nonempty
clause, so minimal members exist and form a finite nonempty family, proving
(c).  If a successful construction has a capability set represented in
\(U\) and the clauses are sound for the exact question, its set escapes; (b)
then supplies the required minimal subset.
\end{proof}

\begin{example}[Two-clause control]
Let \(U=\{d,n,q\}\),
\[
  C_1=\{d,q\},\qquad C_2=\{n,q\}.
\]
Then
\[
  \mathcal H_{\min}=\bigl\{\{q\},\{d,n\}\bigr\}.
\]
The empty set and the singletons \(\{d\}\) and \(\{n\}\) fail.  Here \(d\)
may stand for a source-derived discriminator, \(n\) for naturality and
variation structure, and \(q\) for a quotient with interface descent.  The
labels are capability bookkeeping, not adopted laws.
\end{example}

\begin{proposition}[Incomplete-universe guard]
If a capability \(u_\star\) or a sound discharge mechanism is omitted from
\(U\) or from every relevant clause, Theorem~\ref{thm:hitting} gives no
verdict about extensions using \(u_\star\).  The theorem therefore cannot
establish future-extension impossibility without a separate completeness
result.
\end{proposition}

\begin{proof}
The theorem quantifies only over subsets of \(U\).  An omitted capability is
outside that quantified universe.
\end{proof}

``Minimum'' here always means inclusion-minimal in the declared finite model.
It does not mean physically sufficient, least ontological change, empirically
preferred, or authorized for adoption.

\section{Quotient-first alternatives and their burdens}

Representative choice is not the only mathematical route.  Let
\(I:\C\to\D\) be a functor fixed before any object pair is evaluated and
define
\[
  A\EqI B\quad\Longleftrightarrow\quad I(A)\cong I(B).
\]

\begin{proposition}[Kernel equivalence]
The relation \(\EqI\) is an equivalence relation on the objects of \(\C\).
\end{proposition}

\begin{proof}
Identity, inverse, and composition of isomorphisms in \(\D\) give
reflexivity, symmetry, and transitivity.
\end{proof}

Further properties are conditional:
\begin{enumerate}
  \item a monoidal congruence requires declared operations and natural
        comparison isomorphisms
        \(I(A\otimes B)\cong I(A)\boxtimes I(B)\);
  \item localization preservation requires \(I\) to factor through the
        localization, while the converse requires isomorphism reflection;
  \item exact coarse-graining preservation requires a factorization
        \(I\cong\bar I Q\), while recovering equality before coarse graining
        again requires reflection; and
  \item variation invariance requires a declared natural isomorphism
        \(IV\Rightarrow I\) (or its typed two-category analogue).
\end{enumerate}

The requirement of prospective fixation is necessary for auditability but is
not sufficient for source derivation.  For any partition \(\Pi\) of a finite
set \(X\), discrete
categories and a functor sending each \(x\) to its block realize exactly that
partition as a kernel relation.  The three-element set has five partitions,
including the identity and universal relations.  Thus an independently
derived admissibility or nontriviality principle is still needed.

The marked-interface control is equally important.  Integral graph homology
assigns both a marked edge \(K_2\) and a marked three-vertex path \(P_3\) the
pair \((\beta_0,\beta_1)=(1,0)\).  If an operation adds an edge between the
marked endpoints when absent, it leaves \(K_2\) at \((1,0)\) but turns the
path into a cycle with \((1,1)\).  The homology kernel forgot the interface
fact needed by the operation.  This does not refute the conditional
congruence theorem; it exhibits failure of a missing descent premise.

Consequently, a quotient route needs at least a source-derived functor,
noncircular admissibility, descent for intended interfaces and operations,
adequate preservation/reflection results, and a separate operational bridge.
A mathematical quotient alone is not physical gauge, observational
equivalence, or general \(\mathrm{EqSrc}\).

\section{Relation to established literature}

Eilenberg and Mac Lane's formulation of natural transformations supplies the
coherence vocabulary and the distinction between objectwise and natural
comparisons \cite{eilenberg1945}.  Anonymous-network election results show,
within their own computational hypotheses, how symmetry and added resources
constrain unique choice \cite{angluin1980,itai1981}; they are analogies, not
proofs about the source functor used here.

Palais proves local slice results under scoped Lie-group action hypotheses
\cite{palais1961}, whereas Singer proves a topology-dependent obstruction to a
continuous global gauge choice in a particular gauge-theory setting
\cite{singer1978}.  Neither result licenses a universal EqSrc obstruction.
Marsden and Weinstein exemplify quotient-first reduction under symplectic
reduction hypotheses \cite{marsden1974}; their construction motivates the
alternative route but does not supply this project's source action or
operational equivalence.  Tannaka's reconstruction result illustrates that
canonical reconstruction can require enriched compatible invariant data
\cite{tannaka1939}, not merely a bare orbit set.

Accordingly, the manuscript makes no priority claim for naturality,
fixed-point obstructions, group-action examples, hitting sets, or quotient
constructions individually.  Its novelty claim is limited to the
project-specific integrated presentation and its explicit transfer and
authority guards.  A broader novelty claim would require a dedicated,
independent literature review beyond this internal packet.

\section{Foundations and physical interpretation}

The current canonical Aether-Flow ontology supplies an exact-GR interpretive
benchmark, but it does not derive the category \(\C\), its core \(\G\), the
choice functor \(S\), the transformation \(K\), the extension atoms, or the
invariant functor \(I\) used here.  The mathematical results therefore have
the status \texttt{blocked\_adoption\_open\_continuation}: the present
ontology does not derive the required selector or quotient structures, while
same-milestone research remains open.

Four distinctions must be retained:
\begin{enumerate}
  \item structural isomorphism is not automatically physical gauge;
  \item representative irrelevance is not automatically operational
        equivalence;
  \item a source-extension declaration is not an ontology adoption; and
  \item exact downstream agreement with GR is not a first-principles
        substrate derivation.
\end{enumerate}

No Distance-to-GR ledger row changes.  General \(\mathrm{EqSrc}\), an
unscoped effective metric, matter coupling, source-derived Einstein
equations, benchmark promotion, and a completed derivation remain outside the
proved domain.

\section{Reproducibility and review status}

The task-local proof archive records each theorem, proof method, source hash,
and limitation.  The companion JSON records nine fixed-point examples, all
five partitions of a three-element set, the marked-interface countermodel,
the \(B_8\) counts, and the two minimum signatures of the finite control.
A deterministic checker recomputes these counts and verifies source hashes,
required scope language, and authority flags.  Such validation checks
consistency; it does not create scientific truth or proof authority.

The associated review packet is a same-context internal source-trace and
scope review.  It is not blind independent human review, peer review, or
submission authorization.  External contact, release, and publication remain
human-gated.

\section{Conclusion}

On a declared groupoid and Set-valued choice functor, natural deterministic
selectors are exactly products of automorphism-fixed loci.  This yields
precise obstruction and uniqueness criteria, while a separate relation map
decides whether representative differences survive at the relation level.
Finite examples expose the alternatives and the assumptions that repair them.
Inside a finite sound clause model, every repair contains an
inclusion-minimal hitting signature, but no conclusion follows about omitted
capabilities.  Quotient-first constructions can avoid representative choice
only after their functorial, interface, source-derivation, and operational
burdens are met.

These are scoped mathematical results and an internal integrated manuscript,
not a global rejection of source-equivalence research.  The next scientific
step remains a separately governed plan decision; this manuscript neither
executes it nor supplies protected publication authority.

\begin{thebibliography}{99}

\bibitem{angluin1980}
Angluin, D. (1980).
Local and global properties in networks of processors (Extended abstract).
In \emph{Proceedings of the twelfth annual ACM symposium on Theory of
computing} (pp.~82--93). Association for Computing Machinery.
\url{https://doi.org/10.1145/800141.804655}

\bibitem{eilenberg1945}
Eilenberg, S., \& Mac Lane, S. (1945).
General theory of natural equivalences.
\emph{Transactions of the American Mathematical Society, 58}, 231--294.
\url{https://doi.org/10.1090/S0002-9947-1945-0013131-6}

\bibitem{itai1981}
Itai, A., \& Rodeh, M. (1981).
Symmetry breaking in distributive networks.
In \emph{22nd Annual Symposium on Foundations of Computer Science}
(pp.~150--158). IEEE.
\url{https://doi.org/10.1109/SFCS.1981.41}

\bibitem{marsden1974}
Marsden, J. E., \& Weinstein, A. (1974).
Reduction of symplectic manifolds with symmetry.
\emph{Reports on Mathematical Physics, 5}(1), 121--130.
\url{https://doi.org/10.1016/0034-4877(74)90021-4}

\bibitem{palais1961}
Palais, R. S. (1961).
On the existence of slices for actions of non-compact Lie groups.
\emph{Annals of Mathematics, 73}(2), 295--323.
\url{https://doi.org/10.2307/1970335}

\bibitem{singer1978}
Singer, I. M. (1978).
Some remarks on the Gribov ambiguity.
\emph{Communications in Mathematical Physics, 60}, 7--12.
\url{https://doi.org/10.1007/BF01609471}

\bibitem{tannaka1939}
Tannaka, T. (1939).
\emph{\"Uber den Dualit\"atssatz der nichtkommutativen topologischen Gruppen}.
Tohoku Mathematical Journal, 45, 1--12.
\url{https://www.jstage.jst.go.jp/article/tmj1911/45/0/45_0_1/_article}

\end{thebibliography}

\end{document}
